% ============================================================
\chapter{Stochastic Calculus in Finance}
\label{ch:stochastic_calculus}
% ============================================================

Brownian motion is the fundamental building block of quantitative finance---but it is a mathematically monstrous object: its paths are continuous yet \emph{nowhere differentiable}. One cannot differentiate in the classical sense, and the ordinary integral breaks down. It was Kiyosi It\^o who, in the 1940s in Japan, constructed the differential calculus adapted to this situation: the \emph{It\^o integral} and the celebrated \emph{It\^o formula}, which is to stochastic calculus what the chain rule is to classical calculus---with an additional corrective term, the $\frac{1}{2}\sigma^2 f''$ term, that changes everything. This chapter builds these tools and shows how they apply to financial modelling.

\section{Brownian Motion}

\begin{definition}[Standard Brownian Motion]
A \textbf{standard Brownian motion} (or Wiener process) $(W_t)_{t \geq 0}$
is a continuous stochastic process such that:
\begin{enumerate}
    \item $W_0 = 0$ a.s.
    \item The paths $t \mapsto W_t(\omega)$ are continuous a.s.
    \item Increments are independent: for $0 \leq s < t \leq u < v$,
          $W_t - W_s$ and $W_v - W_u$ are independent.
    \item Increments are Gaussian: $W_t - W_s \sim \mathcal{N}(0, t-s)$.
\end{enumerate}
\end{definition}

\begin{proposition}[Properties of Brownian Motion]
\begin{enumerate}
    \item $\E[W_t] = 0$ and $\Var(W_t) = t$ for all $t \geq 0$.
    \item $\Cov(W_s, W_t) = \min(s,t)$.
    \item $W_t$ is a martingale with respect to its natural filtration.
    \item $W_t^2 - t$ is a martingale.
    \item Paths are a.s.\ nowhere differentiable.
    \item The quadratic variation satisfies $\langle W \rangle_t = t$.
\end{enumerate}
\end{proposition}

\begin{codeblock}[title=\textbf{Simulating Brownian Paths}]
\begin{minted}{python}
import numpy as np
import matplotlib.pyplot as plt

def brownian_paths(T, N, M):
    """Simulate M Brownian paths on [0, T] with N steps."""
    dt = T / N
    dW = np.sqrt(dt) * np.random.standard_normal((M, N))
    W = np.zeros((M, N + 1))
    W[:, 1:] = np.cumsum(dW, axis=1)
    t = np.linspace(0, T, N + 1)
    return t, W

np.random.seed(42)
t, W = brownian_paths(1.0, 1000, 5)
for i in range(5):
    plt.plot(t, W[i])
plt.xlabel('t'); plt.ylabel('$W_t$')
plt.title('Brownian Paths')
plt.grid(True)
plt.show()
\end{minted}
\end{codeblock}

\section{The It\^o Integral}

\subsection{Construction}

The stochastic integral extends the classical integral to random integrands
with respect to Brownian motion.

\begin{definition}[It\^o Integral]
For an adapted process $(H_t)_{t \geq 0}$ satisfying
$\E\!\left[\int_0^T H_t^2\,dt\right] < \infty$, the \textbf{It\^o integral}
is defined as:
\[
    I_T(H) = \int_0^T H_t\,dW_t
           = \lim_{n \to \infty} \sum_{i=0}^{n-1}
             H_{t_i}(W_{t_{i+1}} - W_{t_i}).
\]
The limit is taken in mean square ($L^2$).
\end{definition}

\begin{theorem}[Properties of the It\^o Integral]
\label{thm:ito_integral_props}
Let $I_t = \int_0^t H_s\,dW_s$. Then:
\begin{enumerate}
    \item $I_t$ is a martingale: $\E[I_t \mid \mathcal{F}_s] = I_s$
          for $s \leq t$.
    \item \textbf{It\^o isometry}:
          $\E[I_T^2] = \E\!\left[\int_0^T H_t^2\,dt\right]$.
    \item $\E[I_T] = 0$.
    \item The quadratic variation is
          $\langle I \rangle_t = \int_0^t H_s^2\,ds$.
\end{enumerate}
\end{theorem}

\begin{warning}[title=\textbf{It\^o vs Stratonovich}]
The It\^o integral evaluates the integrand at the \textbf{left endpoint}
of each interval: $H_{t_i}$ rather than $H_{t_{i+1}}$ or
$\frac{1}{2}(H_{t_i} + H_{t_{i+1}})$. The Stratonovich choice (midpoint
evaluation) gives a different result and satisfies classical calculus rules,
but the It\^o integral is preferred in finance because it preserves the
martingale property (and hence risk-neutral pricing).
\end{warning}

\subsection{Calculus Rules}

The rules of calculus differ from the classical case. The key identities are:
\[
    (dW_t)^2 = dt, \quad (dt)^2 = 0, \quad dW_t \cdot dt = 0.
\]

\section{It\^o's Formula}

\begin{theorem}[It\^o's Formula --- One-Dimensional Case]
\label{thm:ito_formula}
Let $(X_t)$ be an It\^o process satisfying
$dX_t = \mu_t\,dt + \sigma_t\,dW_t$
and $f \in C^{1,2}(\R^+ \times \R)$. Then:
\[
    df(t, X_t) = \frac{\partial f}{\partial t}\,dt
    + \frac{\partial f}{\partial x}\,dX_t
    + \frac{1}{2}\frac{\partial^2 f}{\partial x^2}\sigma_t^2\,dt.
\]
In expanded form:
\[
    df(t, X_t) = \left(\frac{\partial f}{\partial t}
    + \mu_t\frac{\partial f}{\partial x}
    + \frac{1}{2}\sigma_t^2\frac{\partial^2 f}{\partial x^2}\right)dt
    + \sigma_t\frac{\partial f}{\partial x}\,dW_t.
\]
\end{theorem}

\begin{keyformulas}[title=\textbf{It\^o's Formula --- Compact Form}]
\[
    df = \left(\partial_t f + \mu\,\partial_x f
    + \tfrac{1}{2}\sigma^2\partial_{xx} f\right)dt
    + \sigma\,\partial_x f\,dW_t.
\]
The extra term $\frac{1}{2}\sigma^2\partial_{xx}f\,dt$ has no analogue
in deterministic calculus; it arises from the non-zero quadratic variation
of Brownian motion.
\end{keyformulas}

\begin{example}
Apply It\^o's formula to $f(x) = x^2$ with $X_t = W_t$:
\[
    d(W_t^2) = 2W_t\,dW_t + \frac{1}{2}\cdot 2\,dt = 2W_t\,dW_t + dt.
\]
Integrating: $W_T^2 = 2\int_0^T W_t\,dW_t + T$, hence:
\[
    \int_0^T W_t\,dW_t = \frac{W_T^2 - T}{2}.
\]
Note the difference from classical calculus, which would give $W_T^2/2$.
\end{example}

\begin{theorem}[Multidimensional It\^o's Formula]
Let $dX_t^i = \mu_t^i\,dt + \sum_j \sigma_t^{ij}\,dW_t^j$ and
$f \in C^{1,2}$. Then:
\[
    df = \partial_t f\,dt + \sum_i \partial_{x_i}f\,dX_t^i
    + \frac{1}{2}\sum_{i,j}\partial_{x_i x_j}f\,
    d\langle X^i, X^j\rangle_t.
\]
\end{theorem}

\section{Girsanov's Theorem}

\begin{theorem}[Girsanov]
\label{thm:girsanov}
Let $(W_t)$ be a $\PP$-Brownian motion and $(\theta_t)$ an adapted process
satisfying the Novikov condition:
\[
    \E^\PP\!\left[\exp\!\left(\frac{1}{2}\int_0^T \theta_t^2\,dt\right)\right]
    < \infty.
\]
Define the Radon--Nikodym density:
\[
    \frac{d\mathbb{Q}}{d\PP}\bigg|_{\mathcal{F}_T}
    = \mathcal{E}_T(\theta) = \exp\!\left(-\int_0^T \theta_t\,dW_t
    - \frac{1}{2}\int_0^T \theta_t^2\,dt\right).
\]
Then under $\mathbb{Q}$, the process
$\tilde{W}_t = W_t + \int_0^t \theta_s\,ds$
is a standard Brownian motion.
\end{theorem}

\begin{intuition}[title=\textbf{Application to Finance}]
Girsanov's theorem enables the passage from the historical measure $\PP$
(under which $dS_t/S_t = \mu\,dt + \sigma\,dW_t$) to the risk-neutral
measure $\mathbb{Q}$ (under which $dS_t/S_t = r\,dt +
\sigma\,d\tilde{W}_t$) by setting $\theta = (\mu - r)/\sigma$.
The quantity $\theta$ is the \textbf{market price of risk}.
\end{intuition}

\section{Martingale Representation Theorem}

\begin{theorem}[Martingale Representation]
\label{thm:representation}
Every square-integrable martingale $(M_t)_{0 \leq t \leq T}$ adapted to
the Brownian filtration can be written as:
\[
    M_t = M_0 + \int_0^t H_s\,dW_s,
\]
for a unique predictable process $(H_t)$ satisfying
$\E\!\left[\int_0^T H_t^2\,dt\right] < \infty$.
\end{theorem}

\begin{remark}
This theorem guarantees that in the Black--Scholes model (complete market),
every square-integrable contingent claim is replicable. The process $H_t$
provides the hedging strategy.
\end{remark}

\section{Stochastic Differential Equations}

\begin{definition}[General SDE]
A \textbf{stochastic differential equation} (SDE) takes the form:
\[
    dX_t = b(t, X_t)\,dt + \sigma(t, X_t)\,dW_t, \quad X_0 = x_0,
\]
where $b$ is the drift coefficient and $\sigma$ the diffusion coefficient.
\end{definition}

\begin{theorem}[Existence and Uniqueness]
If $b$ and $\sigma$ are Lipschitz continuous in $x$ uniformly in $t$ and
have linear growth, then the SDE admits a unique strong solution
$(X_t)_{0 \leq t \leq T}$.
\end{theorem}

\begin{theorem}[PDE--SDE Link (Feynman--Kac)]
\label{thm:feynman_kac}
Let $u(t,x)$ be a solution of:
\[
    \partial_t u + b(t,x)\partial_x u
    + \frac{1}{2}\sigma^2(t,x)\partial_{xx}u - r\,u = 0,
    \quad u(T,x) = g(x).
\]
Then:
\[
    u(t,x) = \E\!\left[e^{-r(T-t)}g(X_T) \mid X_t = x\right],
\]
where $dX_s = b(s,X_s)\,ds + \sigma(s,X_s)\,dW_s$.
\end{theorem}

\begin{remark}
The Feynman--Kac theorem establishes the fundamental link between PDE-based
pricing (Black--Scholes) and risk-neutral expectation pricing (Monte Carlo).
\end{remark}

\section{Stochastic Exponential}

\begin{definition}[Stochastic Exponential]
The \textbf{stochastic exponential} (or Dol\'eans-Dade exponential) of an
It\^o process $X_t$ is:
\[
    \mathcal{E}_t(X) = \exp\!\left(X_t - X_0 - \frac{1}{2}\langle X\rangle_t\right).
\]
It satisfies $d\mathcal{E}_t = \mathcal{E}_t\,dX_t$,
$\mathcal{E}_0 = 1$.
\end{definition}

\begin{proposition}
The GBM $S_t = S_0\exp\!\left[(\mu - \sigma^2/2)t + \sigma W_t\right]$
is the stochastic exponential of $\mu t + \sigma W_t$:
\[
    S_t = S_0\,\mathcal{E}_t(\mu\,\cdot + \sigma W).
\]
\end{proposition}

\section{Implementation: Discretisation Schemes}

\begin{codeblock}[title=\textbf{Euler--Maruyama and Milstein Schemes}]
\begin{minted}{python}
import numpy as np

def euler_maruyama(x0, b, sigma, T, N, M):
    """Euler-Maruyama scheme for dX = b(X)dt + sigma(X)dW."""
    dt = T / N
    X = np.full(M, x0, dtype=float)
    for i in range(N):
        dW = np.sqrt(dt) * np.random.standard_normal(M)
        X = X + b(X)*dt + sigma(X)*dW
    return X

def milstein(x0, b, sigma, sigma_prime, T, N, M):
    """Milstein scheme (strong order 1.0)."""
    dt = T / N
    X = np.full(M, x0, dtype=float)
    for i in range(N):
        dW = np.sqrt(dt) * np.random.standard_normal(M)
        X = (X + b(X)*dt + sigma(X)*dW
             + 0.5*sigma(X)*sigma_prime(X)*(dW**2 - dt))
    return X

# Application to GBM: dS = r*S*dt + sig*S*dW
r_val, sig = 0.05, 0.20
S0, T, N, M = 100, 1.0, 100, 100000

np.random.seed(42)
ST_euler = euler_maruyama(S0, lambda S: r_val*S, lambda S: sig*S,
                          T, N, M)
ST_milstein = milstein(S0, lambda S: r_val*S, lambda S: sig*S,
                       lambda S: sig, T, N, M)
Z = np.random.standard_normal(M)
ST_exact = S0 * np.exp((r_val - 0.5*sig**2)*T + sig*np.sqrt(T)*Z)

print(f"Euler:    E[S_T] = {np.mean(ST_euler):.2f}")
print(f"Milstein: E[S_T] = {np.mean(ST_milstein):.2f}")
print(f"Exact:    E[S_T] = {S0*np.exp(r_val*T):.2f}")
\end{minted}
\end{codeblock}

\section{Exercises}

\begin{exercise}[It\^o's Formula]
\begin{enumerate}
    \item Compute $d(e^{W_t})$ and $d(e^{\alpha W_t - \alpha^2 t/2})$.
    \item Show that $\mathcal{E}_t = e^{\sigma W_t - \sigma^2 t/2}$ is
          a martingale.
    \item Apply It\^o's formula to $f(t,x) = e^{-rt}x$ to verify that
          $e^{-rt}S_t$ is a $\mathbb{Q}$-martingale.
\end{enumerate}
\end{exercise}

\begin{exercise}[Girsanov's Theorem]
\begin{enumerate}
    \item Verify the change of measure explicitly for GBM: write $S_t$
          under $\PP$ then under $\mathbb{Q}$.
    \item Compute the market price of risk $\theta = (\mu-r)/\sigma$
          for $\mu = 0.10$, $r = 0.03$, $\sigma = 0.25$.
    \item Show that the Radon--Nikodym density $\mathcal{E}_T(\theta)$
          is a $\PP$-martingale.
\end{enumerate}
\end{exercise}

\begin{exercise}[Convergence of Schemes]
\begin{enumerate}
    \item Compare the strong error $\E[\abs{S_T^N - S_T}]$ of Euler and
          Milstein schemes as a function of $N$.
    \item Verify that the strong convergence order is $1/2$ for Euler
          and $1$ for Milstein.
    \item Implement a weak order 2 scheme (Platen) and compare.
\end{enumerate}
\end{exercise}

\begin{exercise}[Feynman--Kac]
\begin{enumerate}
    \item Verify the Feynman--Kac theorem for the European call by
          comparing the Black--Scholes PDE solution with the Monte Carlo
          expectation.
    \item Apply Feynman--Kac to the heat equation
          $\partial_t u = \frac{1}{2}\partial_{xx}u$ and interpret.
\end{enumerate}
\end{exercise}
